\chapter{Conclusion and Future Work}\label{ch:conclusion}

\section{Summary of Contributions}

This dissertation built a working adaptive chess-training system and evaluated it rigorously enough to
find out where it does not work as intended. The system analyses real games deterministically, grades
moves by win probability, and measures weakness with a proper denominator (critical positions). It
labels tactics with PuzzleNet, a multi-task neural network trained here on 5.5 million puzzles, which
also estimates how hard a puzzle is and how uncertain that estimate is, and it recommends puzzles with
an online Bayesian IRT learner. The evaluation combined a pre-registered replication, seven paired simulation experiments, a
prequential real-data calibration, and a 300-player real-data study of player modelling. Several
engineering defects that had silently disabled parts of the system were found and documented.

\section{Answers to Research Questions}\label{sec:answers}

\begin{description}
  \item[RQ1 (rating group from play).] \textbf{Yes.} On unseen players, a five-band classifier reaches
        59.1\% accuracy (chance 20\%), is within one band 93\% of the time, and has permutation
        $p=0.002$. No rating information is used. Objective O1 is met for rating group; style was not
        modelled.
  \item[RQ2 (weakness diagnosis from games).] \textbf{Mostly no, at typical history lengths, and not
        because of the labels.} Future per-category misses are only marginally more predictable than
        population base rates, and per-category reliability is near zero. Relabelling all 40{,}067
        critical positions of the 300-player cohort with the neural labeller changes what the positions
        are called --- the two labellers agree on only 37\% of them --- but leaves reliability at zero.
        Shrinking towards population norms remains the best practical approach, and personal weakness has
        to be learned online from puzzle attempts.
  \item[RQ3 (tactic labelling).] \textbf{Yes, once it is learned rather than written.} Hand-written
        rules reach $\kappa = 0.52$ on held-out puzzles, and five of the 24 categories have no rule at
        all. PuzzleNet reaches $\kappa = 0.91$ on the same puzzles and emits every category, and $\kappa =
        0.78$ on engine lines, the condition the application runs in. Downstream, its labels close 84\% of
        the gap between rule-based and perfect labels for early targeting, and what remains is no longer
        statistically distinguishable from perfect. The caveat is that agreement here is agreement with
        Lichess's own automatic tagger, not with a human expert.
  \item[RQ4 (difficulty matching).] \textbf{Yes.} The difficulty-aware learner cuts frustrating puzzles
        from 28--45\% to 1--3\% in every simulated world. On real logs, rating-based pitching is badly
        miscalibrated, and online learning corrects it.
  \item[RQ5 (adaptive vs non-adaptive).] \textbf{Yes, with qualifications.} Adaptive selection beats
        random and round-robin in every simulation, although the original bandit fails two of its
        three pre-registered thresholds. The new policy's benefit comes from difficulty targeting.
        Whether it improves \emph{learning} depends on the learning mechanism, which only a human study
        can establish.
\end{description}

\section{Recommendations for Future Research}\label{sec:future}

The recommendations are ordered by expected value against effort. The first four are realistic before
the submission deadline. The fifth item of the earlier plan, replacing the hand-written tagger with a
learned one, was carried out during this work and is reported in Section~\ref{sec:nn-results}; what
follows is what that component still needs.

\begin{enumerate}
  \item \textbf{Run the human pilot} (1--2 weeks of elapsed time, little effort). Use the pre-registered
        crossover protocol with the existing users to test feasibility. The application already logs
        everything needed.
  \item \textbf{Measure how many games are needed for reliable diagnosis} (hours). The cohort already
        contains 80 games per player, so the reliability curve can be traced as a function of history
        length. This answers the practical question the RQ2 result raises.
  \item \textbf{Discover styles without supervision} (1--2 days). Following
        Jayasekara~\cite{jayasekara2018classification}, cluster the 300 cohort players with PCA and
        Gaussian mixtures on the engine-derived features, and report cluster stability. This addresses
        the style half of Objective O1.
  \item \textbf{Estimate difficulty with a human-move model} (2--3 days, with Maia). Estimate a mined
        puzzle's difficulty as the probability that a Maia model~\cite{mcilroyyoung2020aligning} at the
        player's rating finds the key move, and validate it against Lichess ratings. The same
        probability gives a principled counter-intuitiveness score, which addresses the missing part of
        Objective O2.
  \item \textbf{Give PuzzleNet an architecture that matches the data} (1--2 weeks, needs a GPU). The
        current network is a multilayer perceptron over a hand-built representation, trained on a laptop
        CPU. The published entries for the same public task use convolutional and attention models over
        board tensors~\cite{schutt2024estimating,milosz2024transformers}, which is the obvious next step:
        the encoding and the evaluation harness built here transfer unchanged, so only the trunk needs
        replacing, and the data-scaling curve in Section~\ref{sec:nn-ablations} says whether more capacity
        or more data is the binding constraint.
  \item \textbf{Validate the labels against human judgement} (days, needs expert time). Both taggers are
        scored against Lichess's automatic tagger, so agreement is not accuracy. A few hundred positions
        labelled by a strong player would establish how much of the remaining disagreement is error on
        either side, and would let the harness measure what it claims to measure.
  \item \textbf{Predict difficulty for \emph{this} player, not the population} (1--2 weeks). PuzzleNet
        predicts a Lichess rating, which is a population average. Conditioning the difficulty head on a
        player's rating, or combining it with a human-move model~\cite{mcilroyyoung2020aligning}, would give
        the recommender a per-player difficulty, which is what the targeting rule actually wants.
  \item \textbf{Off-policy evaluation} (once more data exists). With logged propensities,
        inverse-propensity and doubly-robust estimators can score alternative policies from real
        traffic, following Nie et al.~\cite{nie2026discovering}.
  \item \textbf{Reward learning progress, and use context} (weeks). Follow Clement et
        al.~\cite{clement2015multi} in rewarding measured improvement rather than failure, and condition
        on session context.
  \item \textbf{Generative puzzle creation} (beyond an MSc). Either search-based generation (perturb
        mined positions and re-verify with the engine) or the generative RL approach of Feng et
        al.~\cite{feng2025generating}, conditioned on a weakness estimate reliable enough to condition on.
\end{enumerate}

\section{Final Remarks}

The system does what the evidence supports. It adapts difficulty reliably, learns weaknesses online,
and reports its own uncertainty. It does not claim what the evidence refutes, namely a precise diagnosis
of tactical weaknesses from a few dozen games. Several of the most important results are negative, and
some of them are the system's own pre-registered criteria failing. That is not a weakness of the
dissertation: it is what makes its positive claims credible.

\authornote{Add a short personal reflection on the project process, as required by the assessment
criteria (the ``processes'' and ``evaluation'' components of the marking form).}
